\section{E2E}

Os casos E2E derivam principalmente da Seção 9 (fluxos), da Seção 8 (telas e
estados) e da Seção 7 (resultados e regras). O estado atual da interface não é
fonte para inventar comportamento.

Playwright CLI é ferramenta de exploração: snapshots, locators, requests,
console, screenshots e depuração. Playwright Test mantém specs, assertions e
regressão automatizada. A exploração não substitui a suite permanente.

Locators permanentes preferem papel, label e texto semanticamente estável;
\texttt{data-testid} é usado quando necessário. Seletores CSS baseados na
estrutura visual devem ser evitados. Sessões exploratórias independentes não
compartilham estado sem intenção explícita.

O bootstrap Playwright 1.61.1 com browsers Nix está disponível, mas esta rodada
não criou configuração de suite nem teste E2E de negócio. Esses artefatos devem
nascer com o primeiro fluxo real.
